\chapter[Lezione 6 --- 13 ottobre 2026]{Lezione 6: applicazioni, immagini e controimmagini}
\label{chap:lezione-06}

\section{Un esercizio sugli estremi di insiemi separati}
\label{sec:lez06-insiemi-separati}

La separazione introdotta nella \cref{sec:lez03-numeri-reali}
si può esprimere mediante gli estremi della
\cref{sec:lez04-estremi}.

\begin{exercise}[Estremi di insiemi separati]
\label{ex:lez06-estremi-separati}
Siano \(A,B\subseteq\R\) non vuoti e si supponga che
\[
  a\le b\qquad\forall a\in A,\ \forall b\in B.
\]
Dimostrare che \(\sup A\le\inf B\).
\end{exercise}

\begin{proof}[Soluzione dell'esercizio]
Scelti \(a_0\in A\) e \(b_0\in B\), l'ipotesi implica
\[
  a_0\le b\quad\forall b\in B,
  \qquad a\le b_0\quad\forall a\in A.
\]
Per le \cref{def:maggiorante,def:minorante}, entrambi gli
estremi richiesti sono quindi reali.

Ogni \(b\in B\) è un maggiorante di \(A\).
Per la \cref{def:estremo-superiore}, si ha
\[
  \sup A\le b\qquad\forall b\in B.
\]
Ne segue che \(\sup A\) è un minorante di \(B\).
Per la \cref{def:estremo-inferiore}, si conclude
\[
  \sup A\le\inf B.
\]
\end{proof}

\section{Prodotto cartesiano e applicazioni}
\label{sec:lez06-applicazioni}

Per descrivere una corrispondenza tra due insiemi
si considerano le coppie ordinate formate dai
rispettivi elementi.

\begin{definition}[Prodotto cartesiano]
\label{def:lez06-prodotto-cartesiano}
Siano \(A\) e \(B\) due insiemi. Il loro
\emph{prodotto cartesiano} è l'insieme
\[
  A\times B:=\{(a,b):a\in A,\ b\in B\}.
\]
\end{definition}

Un sottoinsieme \(\Phi\subseteq A\times B\)
descrive una corrispondenza tra gli elementi
di \(A\) e quelli di \(B\). Non ogni insieme
di coppie ordinate definisce tuttavia
un'applicazione.

\begin{definition}[Applicazione]
\label{def:lez06-applicazione}
Un insieme \(\Phi\subseteq A\times B\) si dice
\emph{applicazione} da \(A\) a \(B\) se
\[
  \forall a\in A\;\exists!\,b\in B:
  (a,b)\in\Phi.
\]
\end{definition}

Indicando con \(\phi(a)\) l'unico elemento associato
a \(a\in A\), si scrive
\[
  \phi\colon A\longrightarrow B,
  \qquad a\longmapsto\phi(a).
\]
Per ogni \(a\in A\) e \(b\in B\), si ha
\[
  \phi(a)=b
  \quad\Longleftrightarrow\quad
  (a,b)\in\Phi.
\]
L'insieme \(\Phi\) costituisce quindi il
\emph{grafico} dell'applicazione \(\phi\):
\[
  \Phi=\{(a,\phi(a)):a\in A\}.
\]

\begin{figure}[htbp]
  \centering
  \begin{tikzpicture}[unipd/diagram]
  \foreach \shift in {0,4.4,8.8} {
    \begin{scope}[xshift=\shift cm]
      \draw[unipd/edge] (0,0) ellipse (0.55 and 1.15);
      \draw[unipd/edge] (2.35,0) ellipse (0.55 and 1.15);
      \node at (0,1.45) {\(A\)};
      \node at (2.35,1.45) {\(B\)};
      \foreach \y in {-.55,.55} {
        \fill[unipd-primary] (0,\y) circle (2pt);
        \fill[unipd-primary] (2.35,\y) circle (2pt);
      }
    \end{scope}
  }
  \draw[unipd/directed edge] (0,.55) -- (2.35,.55);
  \draw[unipd/directed edge] (4.4,.55) -- (6.75,.55);
  \draw[unipd/directed edge] (4.4,.55) -- (6.75,-.55);
  \draw[unipd/directed edge] (4.4,-.55) -- (6.75,-.55);
  \draw[unipd/directed edge] (8.8,.55) -- (11.15,.55);
  \draw[unipd/directed edge] (8.8,-.55) -- (11.15,.55);
  \node[align=center] at (1.18,-1.55) {valore mancante};
  \node[align=center] at (5.58,-1.55) {due valori};
  \node[align=center] at (9.98,-1.55) {applicazione};
\end{tikzpicture}

  \caption{Solo la corrispondenza a destra è
  un'applicazione; le altre violano rispettivamente
  l'esistenza e l'unicità del valore associato.}
  \label{fig:lez06-relazioni}
\end{figure}

Per una corrispondenza rappresentata nel piano,
il \emph{test della retta verticale} permette
di verificare l'unicità del valore associato:
nessuna retta verticale deve intersecare il
grafico in più di un punto. Nella
\cref{fig:lez06-test-verticale}, la parabola
\(y=x^2\), sul dominio rappresentato \([-1,1]\),
soddisfa tale condizione, mentre la curva
\(x=y^2\) e la circonferenza la violano.
Per definire un'applicazione occorre inoltre
che ogni retta con ascissa nel dominio
intersechi il grafico almeno una volta.

\begin{figure}[htbp]
  \centering
  \begin{tikzpicture}[unipd/diagram,scale=.9]
  \foreach \shift in {0,4.6,9.2} {
    \begin{scope}[xshift=\shift cm]
      \draw[->,unipd-muted] (-1.35,0) -- (1.45,0)
        node[right] {\(x\)};
      \draw[->,unipd-muted] (0,-1.25) -- (0,1.4)
        node[above] {\(y\)};
    \end{scope}
  }
  \draw[unipd/edge,very thick,domain=-1:1,samples=70]
    plot (\x,{\x*\x});
  \draw[unipd-accent,dashed] (.5,-1.15) -- (.5,1.25);
  \fill[unipd-accent] (.5,.25) circle (1.7pt);
  \node at (0,-1.65) {\(y=x^2\)};

  \begin{scope}[xshift=4.6cm]
    \draw[unipd/edge,very thick,domain=-1:1,samples=70]
      plot ({\x*\x},\x);
    \draw[unipd-accent,dashed] (.36,-1.15) -- (.36,1.25);
    \foreach \y in {-.6,.6} {
      \fill[unipd-accent] (.36,\y) circle (1.7pt);
    }
    \node at (0,-1.65) {\(x=y^2\)};
  \end{scope}

  \begin{scope}[xshift=9.2cm]
    \draw[unipd/edge,very thick] (0,0) circle (1);
    \draw[unipd-accent,dashed] (.6,-1.15) -- (.6,1.25);
    \foreach \y in {-.8,.8} {
      \fill[unipd-accent] (.6,\y) circle (1.7pt);
    }
    \node at (0,-1.65) {\(x^2+y^2=1\)};
  \end{scope}
\end{tikzpicture}

  \caption{Test della retta verticale: la prima
  curva presenta un'unica intersezione per ogni
  ascissa del dominio; le altre due ne presentano
  più di una per le ascisse evidenziate.}
  \label{fig:lez06-test-verticale}
\end{figure}

\section{Dominio, codominio e immagine}
\label{sec:lez06-dominio-immagine}

\begin{definition}[Dominio e codominio]
\label{def:lez06-dominio-codominio}
Sia \(f\colon A\longrightarrow B\). L'insieme
\(A\) si dice \emph{dominio} di \(f\) e si indica
con \(D(f)=A\); l'insieme \(B\) si dice
\emph{codominio} e si indica con \(C(f)=B\).
\end{definition}

L'espressione che determina i valori di una
funzione non basta a individuarla: occorre
specificarne anche il dominio e il codominio.
Si considerino, ad esempio, le applicazioni
\[
  \begin{aligned}
    f\colon\R&\longrightarrow\R,
      &x&\longmapsto x^2,\\
    g\colon\R&\longrightarrow[0,+\infty),
      &x&\longmapsto x^2,\\
    h\colon[0,+\infty)&\longrightarrow[0,+\infty),
      &x&\longmapsto x^2.
  \end{aligned}
\]
Le tre applicazioni hanno la stessa espressione
algebrica: \(f\) e \(g\) differiscono per il codominio,
\(g\) e \(h\) per il dominio.

\begin{definition}[Immagine di un'applicazione]
\label{def:lez06-immagine}
L'\emph{immagine} di \(f\colon A\to B\)
è l'insieme
\[
  \operatorname{Im}(f)
  :=\{b\in B:\exists a\in A,\ f(a)=b\}
  =f(A).
\]
\end{definition}

Per la \cref{def:lez06-immagine}, si ha
\[
  \operatorname{Im}(f)\subseteq C(f)=B,
\]
e l'inclusione può essere stretta. Per le tre
applicazioni precedenti,
\[
  \operatorname{Im}(f)
  =\operatorname{Im}(g)
  =\operatorname{Im}(h)=[0,+\infty).
\]
Infatti, il quadrato di un numero reale è
non negativo e, per ogni \(y\ge0\), si ha
\(y=(\sqrt y)^2\). La \cref{fig:lez06-quadrato}
mostra il grafico del quadrato e distingue,
nel caso di \(f\), l'immagine dal codominio.

\begin{figure}[htbp]
  \centering
  \begin{tikzpicture}[unipd/diagram,x=1.25cm,y=.78cm]
  \fill[unipd-accent!12] (-2.2,1) rectangle (2.2,4);
  \draw[unipd-muted,dashed] (-2.2,1) -- (2.2,1);
  \draw[unipd-muted,dashed] (-2.2,4) -- (2.2,4);
  \draw[->,unipd-muted] (-2.5,0) -- (2.75,0) node[right] {\(x\)};
  \draw[->,unipd-muted] (0,-1.05) -- (0,5.05) node[above] {\(y\)};
  \draw[very thick,unipd-primary,domain=-2.2:2.2,samples=90]
    plot (\x,{\x*\x});
  \draw[very thick,unipd-accent] (0,1) -- (0,4);
  \draw[very thick,unipd-accent] (-2,0) -- (-1,0);
  \draw[very thick,unipd-accent] (1,0) -- (2,0);
  \foreach \x in {-2,-1,1,2} {
    \draw[unipd-muted] (\x,.08) -- (\x,-.08)
      node[below] {\(\x\)};
  }
  \foreach \y in {-1,1,4} {
    \draw[unipd-muted] (.07,\y) -- (-.07,\y)
      node[left] {\(\y\)};
  }
  \node[anchor=west] at (2.45,3.8) {\(S=[1,4]\)};
  \node[anchor=west,text=unipd-primary] at (1.05,4.55)
    {\(y=x^2\)};
\end{tikzpicture}

  \caption{Per \(f\colon\R\to\R\), l'immagine
  \([0,+\infty)\) è distinta dal codominio \(\R\).
  All'intervallo \(S=[1,4]\) corrispondono
  le ascisse \([-2,-1]\cup[1,2]\).}
  \label{fig:lez06-quadrato}
\end{figure}

\section{Controimmagine di un sottoinsieme}
\label{sec:lez06-controimmagine}

\begin{definition}[Controimmagine]
\label{def:lez06-controimmagine}
Sia \(f\colon A\longrightarrow B\) e sia
\(S\subseteq B\). La \emph{controimmagine}
di \(S\) mediante \(f\) è l'insieme
\[
  f^{-1}(S):=\{a\in A:f(a)\in S\}.
\]
\end{definition}

\begin{figure}[htbp]
  \centering
  \begin{tikzpicture}[unipd/diagram]
  \fill[unipd-accent!13] (0,.5) ellipse (.49 and .75);
  \fill[unipd-accent!13] (5,.5) ellipse (.49 and .75);
  \draw[unipd/edge] (0,0) ellipse (.95 and 1.65);
  \draw[unipd/edge] (5,0) ellipse (.95 and 1.65);
  \draw[unipd-accent,dashed] (0,.5) ellipse (.49 and .75);
  \draw[unipd-accent,dashed] (5,.5) ellipse (.49 and .75);
  \node at (0,1.92) {\(A\)};
  \node at (5,1.92) {\(B\)};
  \node at (2.5,1.32) {\(f\)};
  \node[text=unipd-accent!80!black] at (0,-1.95)
    {\(f^{-1}(S)\)};
  \node[text=unipd-accent!80!black] at (5,-1.95) {\(S\)};

  \draw[unipd/directed edge] (.09,.91) -- (4.91,.91);
  \draw[unipd/directed edge] (.09,.19) -- (4.91,.19);
  \draw[unipd/directed edge] (.09,-.55) -- (4.91,-.98);
  \draw[unipd/directed edge] (.09,-1.2) -- (4.91,-.98);
  \foreach \y/\n in {.91/1,.19/2,-.55/3,-1.2/4} {
    \fill[unipd-primary] (0,\y) circle (2pt);
    \node[left] at (-.08,\y) {\(a_{\n}\)};
  }
  \foreach \y/\n in {.91/1,.19/2,-.98/3} {
    \fill[unipd-primary] (5,\y) circle (2pt);
    \node[right] at (5.09,\y) {\(b_{\n}\)};
  }
\end{tikzpicture}

  \caption{Per \(S=\{b_1,b_2\}\subseteq B\) si ha
  \(f^{-1}(S)=\{a_1,a_2\}\subseteq A\), anche se
  \(f\) non è iniettiva.}
  \label{fig:lez06-controimmagine-insiemi}
\end{figure}

Si indica con
\[
  \mathcal P(A):=\{S:S\subseteq A\}
\]
l'\emph{insieme delle parti} di \(A\), e
analogamente con \(\mathcal P(B)\) quello di
\(B\). Per ogni \(S\in\mathcal P(B)\), la
\cref{def:lez06-controimmagine} determina
un unico sottoinsieme di \(A\), eventualmente
vuoto. Si ottiene quindi un'applicazione
tra insiemi delle parti:
\[
  f^{-1}\colon\mathcal P(B)\longrightarrow\mathcal P(A),
  \qquad S\longmapsto f^{-1}(S).
\]
In particolare, se \(S\cap\operatorname{Im}(f)
=\varnothing\), allora \(f^{-1}(S)=\varnothing\).

\begin{example}[Controimmagini del quadrato]
\label{ex:lez06-controimmagini-quadrato}
Per l'applicazione \(f\colon\R\to\R\) definita
nella \cref{sec:lez06-dominio-immagine}, si ha
\[
  f^{-1}(\{-1\})=\varnothing,
  \qquad
  f^{-1}(\{0\})=\{0\},
\]
\[
  f^{-1}([1,4])=[-2,-1]\cup[1,2].
\]
La controimmagine dell'intervallo è
rappresentata nella \cref{fig:lez06-quadrato}.

L'equazione \(z^2=-1\) ammette soluzioni
complesse, ma \(z\mapsto z^2\) non definisce
un'applicazione da \(\C\) a \(\R\): ad esempio,
\[
  (1+\imaginaryunit)^2
  =2\imaginaryunit\notin\R.
\]
Considerando invece l'applicazione
\(q\colon\C\to\C\), \(q(z)=z^2\), si ottiene
\[
  q^{-1}(\{-1\})
  =\{\imaginaryunit,-\imaginaryunit\}.
\]
\end{example}

La notazione \(f^{-1}(S)\) indica una
controimmagine anche quando \(f\) non è
iniettiva né biiettiva. Non presuppone
l'esistenza di un'applicazione inversa
\(B\to A\).

\section{Applicazioni suriettive, iniettive e biiettive}
\label{sec:lez06-proprieta-applicazioni}

\begin{definition}[Suriettività]
\label{def:lez06-suriettiva}
Un'applicazione \(f\colon A\longrightarrow B\)
si dice \emph{suriettiva} se
\(\operatorname{Im}(f)=B\), ossia se
\[
  \forall b\in B\;\exists a\in A:
  f(a)=b.
\]
\end{definition}

La suriettività dipende dal codominio, come illustrato
nell'\cref{ex:lez06-quadrato-proprieta}. Per la
\cref{def:lez06-controimmagine}, la controimmagine di
\(\{b\}\) è non vuota per ogni \(b\in B\) se e solo se
\(f\) è suriettiva.

\begin{definition}[Iniettività]
\label{def:lez06-iniettiva}
Un'applicazione \(f\colon A\longrightarrow B\)
si dice \emph{iniettiva} se, per ogni
\(a_1,a_2\in A\), vale
\[
  a_1\ne a_2
  \quad\Longrightarrow\quad
  f(a_1)\ne f(a_2).
\]
\end{definition}

\begin{exercise}[Caratterizzazione dell'iniettività]
\label{ex:lez06-equivalenza-iniettiva}
Dimostrare che \(f\colon A\to B\) è iniettiva
se e solo se, per ogni \(a_1,a_2\in A\),
\[
  f(a_1)=f(a_2)
  \quad\Longrightarrow\quad
  a_1=a_2.
\]
\end{exercise}

\begin{proof}[Soluzione dell'esercizio]
Siano \(a_1,a_2\in A\) arbitrari.

Se \(f\) è iniettiva e
\(f(a_1)=f(a_2)\), non può valere
\(a_1\ne a_2\): altrimenti, per la
\cref{def:lez06-iniettiva}, si avrebbe
\(f(a_1)\ne f(a_2)\). Pertanto \(a_1=a_2\).

Viceversa, si supponga che
\(f(a_1)=f(a_2)\) implichi \(a_1=a_2\)
per ogni coppia di elementi di \(A\).
Se \(a_1\ne a_2\), non può valere
\(f(a_1)=f(a_2)\), perché si avrebbe
\(a_1=a_2\). Quindi \(f(a_1)\ne f(a_2)\)
e \(f\) è iniettiva.
\end{proof}

\begin{definition}[Biiettività]
\label{def:lez06-biettiva}
Un'applicazione \(f\colon A\longrightarrow B\)
si dice \emph{biiettiva} se è sia iniettiva
sia suriettiva.
\end{definition}

Per le \cref{def:lez06-suriettiva,def:lez06-iniettiva},
un'applicazione è biiettiva se e solo se
ogni elemento del codominio possiede uno
e un solo antecedente nel dominio:
\[
  \forall b\in B\;\exists!\,a\in A:
  f(a)=b.
\]
L'esistenza corrisponde alla \cref{def:lez06-suriettiva};
l'unicità segue dall'\cref{ex:lez06-equivalenza-iniettiva}.

\begin{figure}[htbp]
  \centering
  \begin{tikzpicture}[unipd/diagram]
  \foreach \shift in {0,4.4,8.8} {
    \begin{scope}[xshift=\shift cm]
      \draw[unipd/edge] (0,0) ellipse (0.55 and 1.25);
      \draw[unipd/edge] (2.35,0) ellipse (0.55 and 1.25);
      \node at (0,1.5) {\(A\)};
      \node at (2.35,1.5) {\(B\)};
    \end{scope}
  }
  \foreach \y in {-.65,.65} {
    \fill[unipd-primary] (0,\y) circle (2pt);
    \fill[unipd-primary] (8.8,\y) circle (2pt);
    \fill[unipd-primary] (11.15,\y) circle (2pt);
    \fill[unipd-primary] (6.75,\y) circle (2pt);
  }
  \foreach \y in {-.8,0,.8} {
    \fill[unipd-primary] (2.35,\y) circle (2pt);
    \fill[unipd-primary] (4.4,\y) circle (2pt);
  }
  \draw[unipd/directed edge] (0,.65) -- (2.35,.8);
  \draw[unipd/directed edge] (0,-.65) -- (2.35,-.8);
  \draw[unipd/directed edge] (4.4,.8) -- (6.75,.65);
  \draw[unipd/directed edge] (4.4,0) -- (6.75,.65);
  \draw[unipd/directed edge] (4.4,-.8) -- (6.75,-.65);
  \draw[unipd/directed edge] (8.8,.65) -- (11.15,.65);
  \draw[unipd/directed edge] (8.8,-.65) -- (11.15,-.65);
  \node at (1.18,-1.65) {iniettiva};
  \node at (5.58,-1.65) {suriettiva};
  \node at (9.98,-1.65) {biiettiva};
\end{tikzpicture}

  \caption{Da sinistra: applicazione iniettiva
  non suriettiva, suriettiva non iniettiva e
  biiettiva.}
  \label{fig:lez06-proprieta}
\end{figure}

\section{Tre applicazioni definite dal quadrato}
\label{sec:lez06-esempi-quadrato}

Si riprendono le tre applicazioni definite nella
\cref{sec:lez06-dominio-immagine}, distinguendone
iniettività e suriettività.

\begin{example}[Dominio e codominio del quadrato]
\label{ex:lez06-quadrato-proprieta}
Le applicazioni \(f\) e \(g\) non sono
iniettive, poiché \(1\ne-1\) e
\(1^2=(-1)^2\). Inoltre, \(f\) non è
suriettiva su \(\R\), mentre \(g\) è
suriettiva su \([0,+\infty)\). Nessuna
delle due è quindi biiettiva.

Per ogni \(a>0\), entrambi i domini
contengono le due radici reali di \(a\),
per cui
\[
  f^{-1}(\{a\})=g^{-1}(\{a\})
  =\{-\sqrt a,\sqrt a\}.
\]
Per \(a=0\), le controimmagini sono
\(\{0\}\); se \(a<0\), caso possibile
soltanto nel codominio di \(f\), si ha
\(f^{-1}(\{a\})=\varnothing\).

Restringendo il dominio a \([0,+\infty)\),
l'applicazione \(h\) è iniettiva.
Infatti, se \(x_1,x_2\ge0\) e
\(x_1^2=x_2^2\), allora
\[
  (x_1-x_2)(x_1+x_2)=0.
\]
Se \(x_1+x_2=0\), entrambi i numeri
sono nulli; altrimenti, \(x_1-x_2=0\).
In ogni caso, \(x_1=x_2\).

Poiché ogni \(a\ge0\) è immagine di
\(\sqrt a\in[0,+\infty)\), l'applicazione
\(h\) è anche suriettiva e quindi biiettiva.
In particolare, per ogni \(a\ge0\),
\[
  h^{-1}(\{a\})=\{\sqrt a\}.
\]
\end{example}

Per \(h\) esiste quindi anche
l'\emph{applicazione inversa}, distinta
dalla costruzione sulle parti introdotta
nella \cref{sec:lez06-controimmagine}:
\[
  h^{-1}\colon[0,+\infty)\longrightarrow[0,+\infty),
  \qquad y\longmapsto\sqrt y.
\]
L'applicazione inversa associa a ogni
valore il suo unico antecedente.
La scrittura \(h^{-1}(\{a\})\), in cui
l'argomento è un insieme, indica invece
la controimmagine di tale insieme.
